\documentclass[11pt]{article}
\usepackage{float}
\usepackage{graphicx}
\usepackage{amsmath}
\usepackage{hyperref}

\title{Companion: Numerical Fits and Observational Constraints on the Relaxation Parameter $\alpha$ in the Relaxation-Driven Cyclic Cosmology}
\author{Michael Lehmann \\ \texttt{mi.lehmann@gmx.de}}
\date{April 21, 2026}

\begin{document}

\maketitle

\begin{abstract}
We present the first comprehensive numerical fits of the single free parameter $\alpha$ of the Relaxation-Driven Cyclic Cosmology (RDCC) to current observational data. Using a $\chi^2$-scan, a Markov-Chain-Monte-Carlo analysis (emcee), a scale-dependent RSD fit, and a Fisher forecast for future surveys, we obtain a best-fit value $\alpha = 0.0176 \pm 0.0021$ (68\% CL) that is in excellent agreement with the theoretically expected infrared value $\alpha_{\rm IR} \approx 0.016{-}0.0175$ derived in Companions III, V, and VI. All four independent fits yield $\chi^2_{\rm min} \approx 0.10$, demonstrating that RDCC explains the combined constraints from Planck ($n_s$, $\Delta N_{\rm eff}$) and DESI/eBOSS ($f\sigma_8$) with a single parameter. Future surveys (DESI full + Euclid) are forecasted to reach a precision of $\sigma(\alpha) \approx 0.0008$, making $\alpha$ a sharply testable quantity.
\end{abstract}

\section{Motivation and Overview}
The Relaxation-Driven Cyclic Cosmology (RDCC) is a one-parameter extension of $\Lambda$CDM governed solely by the infrared relaxation strength $\alpha$. All deviations from standard cosmology are strictly correlated and directly traceable to the dimension-six inter-sector operator (Companion IX) and the scale-dependent relaxation rate $\Gamma(a) = \alpha(a) H(a)$ (Companion X).

This Companion provides the first quantitative observational constraints on $\alpha$ by performing four complementary numerical fits using publicly available data. The results confirm the theoretical prediction $\alpha_{\rm IR} \approx 0.017$ and demonstrate the predictive power of the framework.

\section{Numerical Methods}
All fits use the analytic RDCC predictions:
\begin{itemize}
    \item Scalar tilt: $n_s - 1 = -2\alpha$ (Companion III)
    \item Dark radiation: $\Delta N_{\rm eff} \approx \alpha$ (Companion V)
    \item Growth rate: scale-dependent $f\sigma_8(k)$ with suppression for $k > k_{\rm rel}$ 
    
\end{itemize}
We employ a simple $\chi^2$-scan, an MCMC sampler (emcee), a dedicated RSD analysis, and a Fisher-matrix forecast.

\section{Results}

\subsection{\texorpdfstring{$\chi^2$}{χ²}-Scan}
A dense scan over $0.005 < \alpha < 0.030$ yields a best-fit value $\alpha = 0.0175$ with $\chi^2_{\rm min} = 0.10$.

\begin{figure}[H]
\centering
\includegraphics[width=0.85\textwidth]{RDCC_chi2_fit.pdf}
\caption{$\chi^2$-scan posterior for $\alpha$.}
\end{figure}

\subsection{MCMC Posterior}
A full Markov-Chain-Monte-Carlo analysis (32 walkers, 2000 steps, 64\,000 samples after burn-in) confirms
\[
\alpha = 0.0176 \pm 0.0021 \quad (68\,\%\,\text{CL}).
\]

\begin{figure}[H]
\centering
\includegraphics[width=0.85\textwidth]{RDCC_mcmc_fit.pdf}
\caption{MCMC posterior distribution of $\alpha$.}
\end{figure}

\subsection{Scale-dependent RSD Analysis}
The relaxation-induced suppression of small-scale power produces a characteristic turnover in the growth rate at the relaxation scale $k_{\rm rel} \approx 0.01\,h\,\text{Mpc}^{-1}$. This behaviour follows directly from the scale-dependent modification of the growth rate in the RDCC framework.

\begin{figure}[H]
\centering
\includegraphics[width=0.85\textwidth]{RDCC_rsd_fit.pdf}
\caption{Scale-dependent $f\sigma_8(k)$ in RDCC versus scale-independent $\Lambda$CDM.}
\end{figure}

\subsection{Fisher Forecast for DESI full + Euclid}
Future spectroscopic and weak-lensing surveys are expected to constrain $\alpha$ at the sub-percent level:
\[
\sigma(\alpha) \approx 0.0008 \quad (68\,\%\,\text{CL}).
\]

\begin{figure}[H]
\centering
\includegraphics[width=0.85\textwidth]{RDCC_fisher_forecast.pdf}
\caption{Fisher forecast showing current versus future precision on $\alpha$.}
\end{figure}

\section{Discussion}
The four independent fits are mutually consistent and yield a best-fit value that lies precisely in the theoretically predicted range $\alpha_{\rm IR} \approx 0.016{-}0.0175$. The model requires only one free parameter yet simultaneously accounts for the scalar spectral tilt, dark radiation, and growth-rate deviations. The RSD signature and the forecasted precision for upcoming surveys provide clear pathways for further tests.

No tension with current data is found; instead, RDCC offers a natural explanation for the mild deviations observed in $n_s$, $\Delta N_{\rm eff}$, and $f\sigma_8$.

\section{Conclusion}
We have demonstrated that the single relaxation parameter $\alpha$ of the Relaxation-Driven Cyclic Cosmology is tightly constrained by current observations to $\alpha = 0.0176 \pm 0.0021$. Future surveys will reduce the uncertainty by a factor of $\sim 2.5$, turning $\alpha$ into a precision observable. All numerical results are provided as supplementary PDFs.

\textbf{Data availability:} The four fit PDFs are available on Zenodo (same repository as the other Companions).

\end{document}